\documentclass[11pt,a4paper]{article}
\usepackage[T1]{fontenc}
\usepackage[utf8]{inputenc}
\usepackage{newtxtext,newtxmath,microtype}
\usepackage[margin=22mm,headheight=14pt]{geometry}
\usepackage{amsmath,booktabs,tabularx,array,caption,float,tikz,pgfplots,siunitx}
\pgfplotsset{compat=1.18}
\usetikzlibrary{arrows.meta,positioning}
\usepackage{xcolor,fancyhdr,lastpage}
\usepackage[hidelinks]{hyperref}
\definecolor{accent}{HTML}{254C58}
\definecolor{blueplot}{HTML}{2166AC}
\definecolor{redplot}{HTML}{B2182B}
\definecolor{greenplot}{HTML}{28734F}
\hypersetup{pdftitle={Section 2 -- Electrodynamic analysis of directly buried NKT cables},pdfauthor={Tiebert Lefebure}}
\pagestyle{fancy}\fancyhf{}
\fancyhead[L]{\small Finite Elements for Electromagnetic Fields}
\fancyhead[R]{\small NKT 33 kV XLPE-AL}
\fancyfoot[L]{\small Section 2: Electrodynamic analysis}
\fancyfoot[R]{\small \thepage\,/\,\pageref{LastPage}}
\renewcommand{\headrulewidth}{0.3pt}
\setlength{\parindent}{0pt}\setlength{\parskip}{5pt}
\setlength{\emergencystretch}{1.5em}
\setlength{\tabcolsep}{5pt}\renewcommand{\arraystretch}{1.15}
\captionsetup{font=small,labelfont=bf,skip=5pt}
\newcolumntype{Y}{>{\raggedright\arraybackslash}X}
\newcommand{\pagetitle}[2]{{\LARGE\bfseries #1}\par{\large #2}\par\medskip}
\newcommand{\subhead}[1]{\par\medskip{\large\bfseries #1}\par\smallskip}
\newcommand{\unitfield}{\ensuremath{\mathrm{kV/mm}}}
\newcommand{\source}[2]{\href{#1}{\textit{#2}}}
\pgfplotsset{report/.style={width=0.95\linewidth,height=5.2cm,grid=major,
 grid style={gray!20},axis line style={gray!70},tick label style={font=\small},
 label style={font=\small},legend style={font=\small,draw=none,fill=white},
 scaled ticks=false,/pgf/number format/1000 sep={\,},
 every axis plot/.append style={line width=1pt},unbounded coords=jump}}
% Generated from validation/summary.json.
\newcommand{\HealthyE}{4.6513}
\newcommand{\AnalyticalE}{4.66777}
\newcommand{\AnalyticalOuterE}{2.50926}
\newcommand{\HealthyC}{0.22407466}
\newcommand{\AnalyticalC}{0.22407448}
\newcommand{\FieldError}{0.352}
\newcommand{\CapError}{0.00008}
\newcommand{\VoidE}{6.1152}
\newcommand{\VoidXLPE}{5.9759}
\newcommand{\VoidChange}{0.212}
\newcommand{\VoidIncrease}{28.48}
\newcommand{\VoidCapChange}{-0.0403}
\newcommand{\ThickerE}{4.2610}
\newcommand{\ThickerC}{0.20547224}
\newcommand{\FieldReduction}{8.39}
\newcommand{\CapReduction}{8.30}
\newcommand{\BelowMax}{6.64}
\newcommand{\DofReduction}{34.71}
\newcommand{\ChargingPeak}{1.89675}
\newcommand{\ChargingRMS}{1.34121}
\newcommand{\ContourError}{0.0996}
\newcommand{\InnerRMS}{3.3006}
\newcommand{\OuterRMS}{1.7743}
\newcommand{\SCFullLoss}{3.1847e-05}
\newcommand{\IdealLoss}{3.6746e-09}

\newcommand{\CurrentRows}{%
Inner semiconductor & 8.903 & \num{0.0337} & \num{1.06e-07} \\
XLPE & 13.291 & \num{3.3e-12} & \num{0.0229} \\
Outer semiconductor & 17.698 & \num{0.017} & \num{5.32e-08} \\
Gas void (defect) & 10.291 & \num{0} & \num{0.0166} \\%
}

% Original question wording: NMES Assignment 2026, Sections 1 and 2.
% Line-wrap hyphenation is removed; mathematical notation is typeset in LaTeX.
\usepackage{tcolorbox}
\newcommand{\assignmentquestion}[1]{%
  \begin{tcolorbox}[colback=accent!4,colframe=accent!65,
    boxrule=0pt,leftrule=2pt,arc=0pt,boxsep=0pt,
    left=7pt,right=7pt,top=5pt,bottom=5pt,
    before skip=0pt,after skip=8pt,fontupper=\small]
  \textbf{Assignment question.} #1
  \end{tcolorbox}%
}
\newcommand{\questionfiles}[1]{%
  \par\smallskip{\footnotesize\color{accent}
  \textbf{Files / results (relative to Source-code).} #1\par}\medskip
}
\newcommand{\QuestionOneOne}{%
Describe your cable: construction, dimensions, materials, rated voltage, rated
current (ampacity) and the laying conditions for which that rating is given
(burial depth, soil thermal resistivity, ambient temperature, arrangement).
Give a table of electromagnetic and thermal properties with a source for each value.}
\newcommand{\QuestionOneTwo}{%
Choose the environment consistently with the datasheet rating conditions,
and justify its properties.}
\newcommand{\QuestionOneThree}{%
State the sheath bonding you assume: single-point bonded (no net sheath current)
or solidly bonded (circulating sheath currents allowed), and how the armour
and pipe are connected. This choice strongly affects the losses in Section 3.}
\newcommand{\QuestionOneFour}{%
State the excitation you impose: phase-to-ground voltage and phase current,
as peak phasors, with the phase sequence (see the conventions in Appendix A).}
\newcommand{\QuestionTwoOne}{%
Define your computational domain, the equations solved and the boundary
conditions. Where do you truncate the domain, and why? Discuss the role of
the grounded screens: which regions actually influence the field in the insulation?}
\newcommand{\QuestionTwoTwo}{%
Can you simplify the geometry of your cable without degrading the precision
(e.g. disregard details in the geometrical description)? Quantify the effect
of each simplification.}
\newcommand{\QuestionTwoThree}{%
How do the results vary with the mesh refinement? Plot the maximum field in
the insulation and the capacitance versus the number of degrees of freedom,
and identify where local refinement is needed.}
\newcommand{\QuestionTwoFour}{%
Define the admissible field strength at the conductor screen and at the core
screen from the datasheet (or, if absent, from a cited reference).
Is the healthy cable within these limits?}
\newcommand{\QuestionTwoFive}{%
Add a defect to the insulation of your cable. State its nature (e.g. gas void,
water tree, protrusion of the semiconducting screen), its properties, size and
position, and discuss what a defect in a 2D cross-section represents in 3D.
Compare the field with and without defect along a radial cut.
Does the cable withstand the defect?}
\newcommand{\QuestionTwoSix}{%
Compute the resistive and displacement current densities in the cable.
Is current conservation satisfied? Which current dominates in each material, and why?}
\newcommand{\QuestionTwoSeven}{%
Compute the per-unit-length capacitance (\(\mu\mathrm{F/km}\)) and compare with
the analytical value (13) and with the datasheet. State the assumptions of the
analytical formula (in particular, how the semiconducting layers are treated).
Does the capacitance vary with frequency? Why (not)?}
\newcommand{\QuestionTwoEight}{%
Propose an improvement of the cable design with respect to the electric stress,
and support it quantitatively.}

\begin{document}

\pagetitle{Domain, equations and boundary conditions}{Q2.1 --- A screened interior electric-field problem}
\assignmentquestion{\QuestionTwoOne}
\questionfiles{\path{section_2/cable.geo} defines the domain; \path{section_2/cable.pro} defines equations/BCs; \path{section_2/settings.inc} controls cases. All Section 2 solves also use \path{section_1/cable_inputs.json}, \path{shared/materials.json}, \path{section_2/prepare.py} and \path{shared/Jacobian_Integration.pro}.}
The baseline is three NKT \(33~\mathrm{kV}\), \(1\times240+50~\mathrm{mm^2}\)
single-core cables, touching in trefoil and directly buried, without armour or
pipe [1,2]. All three phase interiors are solved in their physical positions.
The copper screens are ideal grounded electrodes. This is a uniform 2D
cross-section; joints, terminations, axial voltage drop and helical lay are excluded.

For peak phasors, \(x(t)=\Re\{\widehat x e^{\mathrm j\omega t}\}\), solve [1]
\[
 -\nabla\!\cdot\!\big[(\sigma+\mathrm j\omega\varepsilon)\nabla\widehat v\big]=0,
 \qquad \widehat{\mathbf E}=-\nabla\widehat v,\qquad \omega=2\pi(50)~\mathrm{rad/s}.
\]
The weak form is \(\int_\Omega(\sigma+\mathrm j\omega\varepsilon)
\nabla\widehat v\cdot\nabla w\,\mathrm dA=0\), for test functions vanishing on
the electrode boundaries. It uses the TA's resistive term and displacement-current
term, with first-order nodal triangles and the supplied quadrature [1,8].

\begin{table}[H]\centering\small
\caption{Solved regions in each phase; \(\varepsilon=\varepsilon_0\varepsilon_r\).}
\begin{tabularx}{\linewidth}{@{}Y l c c@{}}
\toprule
Region & Radius interval [mm] & \(\varepsilon_r\) & \(\sigma\) [S/m]\\
\midrule
Inner semiconductor & 8.500--9.300 & 2.25 & \(0.002\)\\
XLPE & 9.300--17.300 & 2.5 & \(10^{-18}\)\\
Outer semiconductor & 17.300--18.100 & 2.25 & \(0.002\)\\
Equivalent contact bedding & 18.100--19.751 & 2.25 & \(0.002\)\\
\bottomrule
\end{tabularx}
\end{table}
The semiconductor and contact-bedding properties are the Section 1 reference
compound [3,4]; XLPE properties follow [5,6]. The contact region is a geometric
homogenisation, not a claimed extra NKT extrusion.

\subhead{Physical electrodes define the truncation}
At the aluminium surface \(r=8.5~\mathrm{mm}\), impose the phase voltage:
\[
 \widehat V_k=\sqrt{\frac23}(33~\mathrm{kV})e^{-\mathrm j2\pi k/3},
 \quad k=0,1,2\;\text{for A, B, C}.
\]
The amplitude is \(26.9444~\mathrm{kV}\); A is above B and C. At the inner
copper-screen boundary, \(r=19.7511~\mathrm{mm}\), impose \(\widehat v=0\).
Material interfaces share nodes: potential is continuous and normal total
current is continuous in the weak formulation. There is no artificial far-field
boundary inside the solved dielectric domain.

\textbf{Conductor interior.} An ideal equipotential conductor has
\(\nabla\widehat v=0\). Its voltage is represented by its boundary, so the
unmeshed central disk is aluminium, not a physical gas void.

\textbf{Jacket and soil.} A continuous grounded screen isolates the interior.
With all outer screen surfaces and remote earth at zero, the passive exterior
problem has \(\widehat v=0\) and \(\widehat{\mathbf E}=0\). Thus jacket, soil
and neighbouring phases do not affect the insulation field \emph{within this
ideal-screen model}. The wire-screen approximation is assessed in Q2.2.

\textbf{Scope of excitation.} The \(415~\mathrm A\) longitudinal load is reserved
for Sections 3--4. Here the currents are transverse charging/leakage currents.
Quoted field norms are phasor norms; for the healthy radial sinusoidal field,
they equal instantaneous peak values. RMS field values are smaller by \(\sqrt2\).

\newpage
\pagetitle{Simplifications and their consequences}{Q2.2 --- Which details can be removed?}
\assignmentquestion{\QuestionTwoTwo}
\questionfiles{Common electric model (Q2.1), plus \path{section_2/study.py} and \path{section_2/analyse.py}. Compare \path{section_2/results/healthy/metrics.json}, \path{section_2/results/ideal/metrics.json} and \path{section_2/validation/analytical_reference.json}.}
Simplifications are judged against a stated quantity. A choice that is adequate
for insulation field and capacitance need not preserve losses, magnetic fields
or temperature. The reference FE model retains the finite-conductivity
semiconductors and equivalent contact bedding.

\begin{table}[H]\centering\small
\caption{Model reductions and the evidence supporting them.}
\begin{tabularx}{\linewidth}{@{}p{36mm} Y Y@{}}
\toprule
Choice & Effect on this electric problem & Limit of the claim\\
\midrule
Omit aluminium interior & Zero change under an ideal prescribed-voltage electrode; interior field and energy vanish. & Axial resistance, skin effect and conductor loss require a conductor mesh in Section 3.\\
Omit copper interior, PE and soil & Zero exterior field/energy under a continuous screen at zero potential. No distant-boundary truncation error is introduced. & No conclusion about real screen gaps, ground-potential rise or exterior magnetic/thermal fields.\\
Continuous copper annulus & Replaces discrete wires/tape with a closed equipotential boundary at the established inner radius. & Real wire/tape shielding error is not quantified: detailed lay and tape geometry are unavailable.\\
Uniform 2D cross-section & Represents a long, straight section without axial variation. & End effects and helical details are not validated by the coaxial check.\\
Omit semiconductors from the mesh & Impose phase/earth voltages directly on the XLPE interfaces at 9.3/17.3 mm. Quantified below. & Preserves field/C here, but removes semiconductor conduction losses.\\
\bottomrule
\end{tabularx}
\end{table}

\subhead{Quantified comparison: finite versus ideal semiconductors}
The ideal-screen model uses \(36\,761\) complex unknowns instead of
\(56\,305\): \textbf{\DofReduction\% fewer}. At the fine mesh, changes in
phase-A capacitance and maximum XLPE field are below \(10^{-8}\) relative.
This does not establish eight-digit accuracy of the FE field: its discretisation
error is much larger (Q2.3).

An independent cylindrical multilayer solution separates the physical reduction
from mesh effects. For each layer,
\[
 Z'_\ell=\frac{\ln(r_{o,\ell}/r_{i,\ell})}
 {2\pi(\sigma_\ell+\mathrm j\omega\varepsilon_\ell)},\qquad
 Y'=\frac{1}{\sum_\ell Z'_\ell},\qquad C'=\frac{\Im Y'}{\omega}.
\]
At 50 Hz the ideal-screen approximation changes the analytical capacitance by
about \(5.5\times10^{-12}\) relative and the inner XLPE field amplitude by
\(4.7\times10^{-12}\). These tiny values are consequences of the assumed
constant semiconductor conductivity, not measurements of the real cable.

\textbf{Losses are not preserved.} The phase-A time-average transverse conductive
loss is approximately \(3.18\times10^{-5}~\mathrm{W/m}\) with finite screens,
versus \(3.67\times10^{-9}~\mathrm{W/m}\) for XLPE alone. Idealisation removes
almost all of this small contribution. It also omits dielectric polarisation loss
because no loss tangent is included. These are not the total operating cable losses.

\textbf{Decision.} Retain finite semiconductors in the documented baseline;
use the XLPE-only model as a verified field/capacitance simplification.
Do not carry these omissions unchanged into the magnetic or thermal models.

\newpage
\pagetitle{Mesh convergence}{Q2.3 --- Field maxima converge more slowly than capacitance}
\assignmentquestion{\QuestionTwoThree}
\questionfiles{\path{section_2/study.py} varies mesh size in \path{section_2/cable.geo}; \path{section_2/analyse.py} extracts maxima/C. Outputs: \path{section_2/validation/convergence.csv}, \path{field_convergence.pos} and \path{capacitance_convergence.pos} in that validation folder.}
The nominal bulk element size is reduced from 1 to 0.5 to 0.25 mm.
Circular CAD interfaces are approximated by straight mesh edges, so boundary
geometry and the P1 potential approximation improve together. The table reports
phase A; the fine-mesh capacitance spread between phases is only
\(1.86\times10^{-5}\%\), consistent with electrically independent screened phases.

\begin{table}[H]\centering\small
\caption{Healthy-cable refinement; all three phase interiors are included.}
\begin{tabular}{@{}r r r r r@{}}
\toprule
\(h\) [mm] & Complex DOFs & Triangles & \(C'\) [\(\mu\mathrm F/\mathrm{km}\)] & \(E_{\max}\) [peak \unitfield]\\
\midrule
1.00 & 3\,645 & 7\,842 & 0.22407752 & 4.57877\\
0.50 & 14\,157 & 29\,394 & 0.22407594 & 4.62530\\
0.25 & 56\,305 & 114\,758 & \HealthyC & \HealthyE\\
\bottomrule
\end{tabular}
\end{table}
\begin{figure}[H]\centering
\begin{tikzpicture}
\begin{axis}[report,height=4.3cm,xlabel={Complex degrees of freedom},ylabel={Peak XLPE field [kV/mm]},
 xtick={3645,14157,56305},xticklabels={3\,645,14\,157,56\,305},ymin=4.56,ymax=4.69,
 legend pos=south east]
\addplot[blueplot,mark=*] table[x=dof,y=field,col sep=comma]{convergence.csv};
\addlegendentry{FE maximum}
\addplot[redplot,dashed,domain=3645:56305]{4.66776671726};\addlegendentry{Analytical inner interface}
\end{axis}
\end{tikzpicture}
\par\smallskip
\begin{tikzpicture}
\begin{axis}[report,height=4.3cm,xlabel={Complex degrees of freedom},ylabel={\(C'\) [\(\mu\mathrm F/\mathrm{km}\)]},
 xtick={3645,14157,56305},xticklabels={3\,645,14\,157,56\,305},
 ymin=0.224074,ymax=0.2240785,ytick={0.224074,0.224076,0.224078},
 yticklabel style={/pgf/number format/fixed,/pgf/number format/precision=6},legend pos=north east]
\addplot[blueplot,mark=*] table[x=dof,y=cap,col sep=comma]{convergence.csv};\addlegendentry{FE phase A}
\addplot[redplot,dashed,domain=3645:56305]{0.22407447529896};\addlegendentry{Analytical coaxial}
\end{axis}
\end{tikzpicture}
\caption{Independent convergence of a local maximum and an integrated quantity.}
\end{figure}
The analytical values are \(\AnalyticalE~\unitfield\) and
\(\AnalyticalC~\mu\mathrm F/\mathrm{km}\). Fine-mesh errors are
\textbf{\FieldError\% in the maximum field} and \textbf{\CapError\% in capacitance}.
The last field increment is about 0.56\%; three levels demonstrate convergence,
without proving an asymptotic error bound.

Refine near the inner XLPE interface, material boundaries and the void boundary.
Raw P1 gradients are constant per triangle and are not smoothed for the reported
maxima. The tested real-arithmetic GetDP build has twice as many real unknowns
as the complex DOF count shown here. Numerical tolerances are not insulation limits.

\newpage
\pagetitle{Admissible electric stress}{Q2.4 --- A sourced comparison, with an explicit applicability limit}
\assignmentquestion{\QuestionTwoFour}
\questionfiles{\path{section_2/analyse.py} exports interface fields to \path{section_2/results/healthy/metrics.json}. Compare them with the cited limits below. Source/limit interpretation is documented in \path{section_2/README.md}; it is not a solver boundary condition.}
The selected NKT sheet specifies voltage and insulation dimensions, but does
not give admissible local fields [2]. Adopt external comparison values from
Eskom specification 240-56030625, Rev.~3, \S3.3(i) and Annex B, items 3.29--3.30
[7]: \(8~\unitfield\) at the conductor screen and \(4~\unitfield\) at the
insulation/core screen, at \(U_0\). Since \(U_0\) is RMS, compare RMS fields.

\textbf{Applicability.} That specification concerns 44--132 kV XLPE cables.
Using its values for this 33 kV example is an explicitly adopted engineering
comparison, not an NKT acceptance criterion. Neither value is a universal XLPE
breakdown strength.

\begin{table}[H]\centering\small
\caption{Healthy stresses on the XLPE side of each semiconducting interface.}
\begin{tabularx}{\linewidth}{@{}Y r r r@{}}
\toprule
Location & FE [RMS \unitfield] & Analytical [RMS \unitfield] & Comparison limit\\
\midrule
Inner XLPE, \(r_i=9.3~\mathrm{mm}\) & 3.2614 & \InnerRMS & 8.0\\
Outer XLPE, \(r_o=17.3~\mathrm{mm}\) & 1.7920 & \OuterRMS & 4.0\\
\bottomrule
\end{tabularx}
\end{table}
The FE interface entries use maxima from triangles with an edge on that
interface. They differ slightly from the global FE maximum and from exact
boundary values. In the healthy coaxial problem, use the analytical boundary
values for the clearest comparison:
\[
 E_{i,\mathrm{rms}}=\frac{33/\sqrt3}{r_i\ln(r_o/r_i)},\qquad
 E_{o,\mathrm{rms}}=\frac{33/\sqrt3}{r_o\ln(r_o/r_i)},
\]
with voltage in kV and radii in mm. Their peak counterparts are
\(\AnalyticalE\) and \(\AnalyticalOuterE~\unitfield\).

\subhead{What follows from the calculation}
\textbf{The healthy model is below both adopted comparison limits.}
The analytical RMS stresses are approximately 41\% and 44\% of those limits.
The remaining numerical boundary-field errors are about 1.2\% at the inner
interface and 1.0\% at the outer interface, too small to alter this comparison.

This assessment applies to smooth, concentric insulation and the chosen material
constants. It does not establish manufacturer-specific withstand, ageing life,
impulse performance or partial-discharge inception. Surface protrusions, voids,
contamination and local contact defects are excluded from the healthy geometry.

\subhead{Why the interfaces matter}
The \emph{conductor screen} is the inner semiconductor; its outer surface is the
inner XLPE boundary. The \emph{core/insulation screen} is the outer semiconductor;
its inner surface is the outer XLPE boundary. The high stresses are evaluated
\textbf{in XLPE adjacent to these surfaces}, not inside the almost equipotential
semiconducting material or at the copper-metal radius.

The maximum in a healthy coaxial annulus occurs at its smallest insulation
radius because \(E\propto1/r\). This explains the inner-interface refinement in
Q2.3 and motivates the insulation-thickness study in Q2.8.

\newpage
\pagetitle{A controlled insulation defect}{Q2.5 --- An axial gas channel in phase A}
\assignmentquestion{\QuestionTwoFive}
\questionfiles{\path{section_2/settings.inc} sets the void; \path{section_2/cable.geo} meshes it; \path{section_2/cable.pro} assigns gas properties. Use \path{section_2/study.py}; inspect \path{section_2/results/void/view.geo}, \path{radial_cut.csv} in that case, and \path{section_2/validation/radial_comparison.pos}.}
Introduce a circular gas region of radius \(0.25~\mathrm{mm}\), centred
\(10.3~\mathrm{mm}\) from the phase-A axis on the positive local \(x\)-direction.
Its nearest edge is \(0.75~\mathrm{mm}\) from the inner XLPE interface.
The position and size are investigation choices, not manufacturing tolerances.
Set \(\varepsilon_r=1\), \(\sigma=0\): an ideal dry gas before discharge.
No gas pressure, discharge threshold or ionisation law is supplied to the solver.

In 2D this disk represents an \textbf{infinitely long axial channel}.
It does not represent a spherical bubble; a finite bubble has different 3D
polarisation and end effects. The channel model is useful for a controlled
field-concentration study, not a prediction of an actual NKT defect.

\begin{table}[H]\centering\small
\caption{Defect results and local refinement check (phase A).}
\begin{tabularx}{\linewidth}{@{}Y r r@{}}
\toprule
Quantity & Coarser void mesh & Fine void mesh\\
\midrule
Complex DOFs (all phases) & 15\,336 & 61\,106\\
Maximum gas field [peak \unitfield] & 6.1023 & \VoidE\\
Maximum XLPE field [peak \unitfield] & 5.9693 & \VoidXLPE\\
Phase-A capacitance [\(\mu\mathrm F/\mathrm{km}\)] & 0.22398569 & 0.22398426\\
\bottomrule
\end{tabularx}
\end{table}
The gas maximum changes by \VoidChange\% on the last refinement. The gas boundary
uses a target size of \(h/12\), about \(0.0208~\mathrm{mm}\) in the fine case,
grading back to the bulk size over 2 mm. The maximum XLPE stress increases by
\VoidIncrease\% relative to the healthy FE maximum, whereas capacitance changes
by only \VoidCapChange\%: a small global change can conceal a local defect.

\begin{figure}[H]\centering
\begin{tikzpicture}
\begin{axis}[report,height=4.3cm,xmin=9.3,xmax=12,ymin=2.2,ymax=6.5,
 xlabel={Radius from phase-A axis [mm]},ylabel={Phasor field norm [kV/mm]},legend pos=north east]
\addplot[blueplot] table[x=radius,y=field,col sep=comma]{healthy_cut.csv};\addlegendentry{Healthy}
\addplot[redplot] table[x=radius,y=field,col sep=comma]{void_cut.csv};\addlegendentry{Gas channel}
\draw[densely dashed,gray] (axis cs:10.05,2.2)--(axis cs:10.05,6.5);
\draw[densely dashed,gray] (axis cs:10.55,2.2)--(axis cs:10.55,6.5);
\end{axis}
\end{tikzpicture}
\caption{Positive radial cut through the channel. Dashed lines mark its edges.
The plotted cut need not pass through the global field maximum.}
\end{figure}
The field jumps at the gas/XLPE interface because permittivity changes. Normal
total current, not electric field, is continuous. Gas stress increases while
nearby XLPE stress redistributes; tangential field remains continuous in the
continuum solution.

\textbf{Withstand is not established.} The linear model contains no discharge,
space charge or ageing. The XLPE comparison limits of Q2.4 cannot be applied
to gas. Additional pressure-dependent inception data, a discharge model or
appropriate tests are required before accepting this defective cable.

\newpage
\pagetitle{Current densities and conservation}{Q2.6 --- Transverse conduction and displacement currents}
\assignmentquestion{\QuestionTwoSix}
\questionfiles{\path{section_2/cable.pro} exports current-density maps; \path{section_2/analyse.py} checks conservation. Outputs in \path{section_2/results/healthy/}: \path{view.geo}, \path{j_conduction.pos}, \path{j_displacement.pos}, \path{radial_cut.csv} and \path{metrics.json}.}
The solved field gives the peak phasors
\[
 \widehat{\mathbf J}_{c}=\sigma\widehat{\mathbf E},\qquad
 \widehat{\mathbf J}_{d}=\mathrm j\omega\varepsilon\widehat{\mathbf E},\qquad
 \widehat{\mathbf J}_{\mathrm{tot}}=\widehat{\mathbf J}_{c}+\widehat{\mathbf J}_{d}.
\]
These currents cross the insulation and semiconducting layers. They are not
the longitudinal \(415~\mathrm A\) load current. In a uniform material their
magnitude ratio is independent of field strength:
\[
 \frac{\|\widehat{\mathbf J}_c\|}{\|\widehat{\mathbf J}_d\|}
 =\frac{\sigma}{\omega\varepsilon}.
\]
At 50 Hz this is \(3.20\times10^5\) in the semiconductor/contact compound,
\(1.44\times10^{-10}\) in XLPE, and zero in the ideal gas. Thus conduction
dominates in the screens, displacement current in XLPE and gas. The total
normal current can be continuous even when its dominant mechanism changes.

\begin{table}[H]\centering\small
\caption{Representative computed samples on the positive phase-A radial cut;
current-density magnitudes are peak norms in A/m\(^2\).}
\begin{tabularx}{\linewidth}{@{}Y r r r@{}}
\toprule
Region & Radius [mm] & \(\|\widehat{\mathbf J}_c\|\) & \(\|\widehat{\mathbf J}_d\|\)\\
\midrule
\CurrentRows
\bottomrule
\end{tabularx}
\end{table}
Samples are from the fine healthy case except the gas row, which uses the fine
defect case. They are local values, not maxima. Full radial data and complex
field components are saved in each case's \path{radial_cut.csv}; both magnitude
maps can be displayed in Gmsh.

\subhead{Two complementary conservation checks}
\textbf{Discrete weak conservation.} Post-processing reconstructs the P1 gradients
from GetDP nodal voltages and assembles element reactions \(K_e\widehat v_e\).
The sum at the inner electrode is its injected transverse current per metre;
the outer-electrode sum has opposite sign. For phase A:
\[
 \widehat I'_A\simeq 2.36\times10^{-9}+\mathrm j\,1.89675\times10^{-3}~\mathrm{A/m}.
\]
Its norm is \(\ChargingPeak~\mathrm{A/km}\) peak, or
\(\ChargingRMS~\mathrm{A/km}\) RMS. Terminal imbalance relative to its norm is
\(1.25\times10^{-10}\); the largest free-node weak residual, with the same
normalisation, is \(1.28\times10^{-10}\).

\textbf{Independent contour check.} Integrate the raw total-current field around
circles at radii 10.8, 13.3 and 16.8 mm. The largest difference from the terminal
reaction current is \textbf{\ContourError\%}. This is larger than the weak
residual because raw P1 gradients do not enforce pointwise normal-flux
continuity across every triangle edge. It quantifies a field-discretisation
error rather than a physical current source.

\textbf{Conclusion.} Current conservation is satisfied to the reported weak
residual and contour accuracy. The small in-phase terminal current is sensitive
to cancellation; do not assign high precision to the tiny leakage/loss estimate.
No dielectric-loss tangent is included, so real dielectric heating is not fully represented.

\newpage
\pagetitle{Capacitance and frequency}{Q2.7 --- Consistent terminal, energy and analytical values}
\assignmentquestion{\QuestionTwoSeven}
\questionfiles{\path{section_2/analyse.py} computes terminal/energy C; \path{section_2/study.py} compares frequency cases and the coaxial formula. Read \path{section_2/validation/summary.json}, \path{analytical_reference.json} in that folder, and the NKT datasheet in the repository root.}
For each phase separately, compute the terminal admittance and capacitance:
\[
 Y'=\frac{\widehat I'}{\widehat V},\qquad
 C'_{\mathrm{FE}}=\frac{\Im Y'}{\omega}.
\]
An independent energy integral gives the same effective capacitance for this
linear two-terminal domain:
\[
 C'_{\mathrm{energy}}=
 \frac{\displaystyle\int_{\Omega_k}\varepsilon
       \|\widehat{\mathbf E}\|^2\,\mathrm dA}{|\widehat V_k|^2},
 \qquad
 W'_{\mathrm{avg},k}=\frac14 C'|\widehat V_k|^2.
\]
For peak phasors, the time-average energy uses \(1/4\), and time-average Joule
loss uses \(1/2\). If total three-phase energy were used, its conversion would
also require division by three. The implementation avoids that ambiguity by
integrating each phase separately.

\begin{table}[H]\centering\small
\caption{Phase-A capacitance at 50 Hz. Multiply F/m by \(10^9\) for \(\mu\)F/km.}
\begin{tabularx}{\linewidth}{@{}Y r@{}}
\toprule
Method / reference & \(C'\) [\(\mu\mathrm F/\mathrm{km}\)]\\
\midrule
FE terminal susceptance & \HealthyC\\
FE energy integral & \HealthyC\\
Analytical coaxial, ideal semiconductors & \AnalyticalC\\
NKT datasheet maximum [2] & 0.24000000\\
\bottomrule
\end{tabularx}
\end{table}
The FE energy and terminal values agree to better than \(10^{-10}\) relative
in the healthy case. The \textbf{analytical insulation radii} are
\(r_i=9.3~\mathrm{mm}\), \(r_o=17.3~\mathrm{mm}\), giving [1]
\[
 C'=\frac{2\pi\varepsilon_0\varepsilon_{r,\mathrm{XLPE}}}{\ln(r_o/r_i)}.
\]
This assumes long concentric cylinders, uniform linear XLPE, ideal equipotential
semiconductors, no end fringing and no defect. Using the metal radii instead
would incorrectly treat semiconductors/contact bedding as high-voltage insulation.
The finite-screen analytical correction in Q2.2 is negligible for C at 50 Hz.

The FE result is \BelowMax\% below the datasheet's \emph{maximum}, consistent
with that bound. It is not validation against a measured nominal capacitance.
Unpublished dimensions and the assumed XLPE permittivity prevent a stronger claim;
no parameter is tuned solely to force agreement with the maximum.

\begin{table}[H]\centering\small
\caption{Frequency study with fixed geometry and constant material properties.}
\begin{tabular}{@{}l r r r r@{}}
\toprule
Frequency [Hz] & 5 & 50 & 500 & 5\,000\\
\midrule
FE \(C'\) [\(\mu\mathrm F/\mathrm{km}\)] & 0.224074656 & 0.224074656 & 0.224074656 & 0.224074644\\
\bottomrule
\end{tabular}
\end{table}
For ideal screens and constant \(\varepsilon\), the formula has no frequency
term. Finite screen resistance introduces a very weak dependence (about
\(5.5\times10^{-8}\) relative over the tested range). This study does not
model dielectric dispersion or establish broadband validity of the 2D approximation.

\newpage
\pagetitle{A quantified design improvement}{Q2.8 --- Increase XLPE thickness from 8 to 9 mm}
\assignmentquestion{\QuestionTwoEight}
\questionfiles{\path{section_2/settings.inc} and \path{section_2/cable.geo} define added XLPE; \path{section_2/study.py} compares cases. Inspect \path{section_2/results/thicker_xlpe/metrics.json}, its \path{view.geo}, and \path{section_2/validation/radial_comparison.pos}.}
Keep the aluminium conductor and inner semiconductor unchanged, so
\(r_i=9.3~\mathrm{mm}\). Increase the outer XLPE radius from 17.3 to 18.3 mm.
Retain the outer semiconductor/contact thicknesses and shift the outside
interfaces outwards by 1 mm. This is a controlled design variant, not a new
manufacturer specification.

The analytical reason is direct:
\[
 E_{i,\mathrm{pk}}=\frac{V_{\mathrm{pk}}}{r_i\ln(r_o/r_i)},\qquad
 C'=\frac{2\pi\varepsilon}{\ln(r_o/r_i)}.
\]
At fixed \(r_i\), raising \(r_o\) increases the logarithmic denominator, reducing
both inner electric stress and capacitance. The improvement is checked with
the same nominal \(0.25~\mathrm{mm}\) bulk mesh size and a healthy insulation.

\begin{table}[H]\centering\small
\caption{Baseline versus the hypothetical thicker-insulation variant (phase A).}
\begin{tabularx}{\linewidth}{@{}Y r r@{}}
\toprule
Quantity & NKT-based baseline & Design variant\\
\midrule
XLPE thickness [mm] & 8.0 & 9.0\\
Nominal outside cable diameter [mm] & 45.5 & 47.5\\
Touching-trefoil centre spacing [mm] & 45.5 & 47.5\\
Maximum FE XLPE field [peak \unitfield] & \HealthyE & \ThickerE\\
FE capacitance [\(\mu\mathrm F/\mathrm{km}\)] & \HealthyC & \ThickerC\\
Complex DOFs (all phases) & 56\,305 & 63\,451\\
\bottomrule
\end{tabularx}
\end{table}
The computed maximum field decreases by \textbf{\FieldReduction\%}; capacitance
falls by \textbf{\CapReduction\%}. Both effects are much larger than the healthy
mesh errors, and the analytical formula predicts the same trend.

\begin{figure}[H]\centering
\begin{tikzpicture}
\begin{axis}[report,height=4.3cm,xmin=9.3,xmax=18.5,ymin=0,ymax=5,
 xlabel={Radius from phase-A axis [mm]},ylabel={Phasor field norm [kV/mm]},legend pos=north east]
\addplot[blueplot] table[x=radius,y=field,col sep=comma]{healthy_cut.csv};\addlegendentry{8 mm XLPE}
\addplot[greenplot] table[x=radius,y=field,col sep=comma]{thicker_xlpe_cut.csv};\addlegendentry{9 mm XLPE}
\end{axis}
\end{tikzpicture}
\caption{Field reduction and the outward shift of the outer XLPE interface.}
\end{figure}
\textbf{Trade-offs and geometry bookkeeping.} More insulation increases material
use, outside diameter and thermal resistance. In this variant the copper
annulus thickness is retained while its radius shifts outwards, so its area
would rise from 50 to about \(52.51~\mathrm{mm^2}\). Copper is not meshed in this
electric model, but a manufacturable redesign must specify its area and the
remaining layers consistently.

\textbf{Decision.} The variant is an effective electric-stress improvement.
Its cost, manufacturability and permissible current remain to be assessed.
The original \(415~\mathrm A\) rating is not transferred to the altered cable.
The chosen baseline remains the original 8 mm NKT insulation.

\newpage
\pagetitle{Sources and reproducibility}{References and implementation record}
Numerical values come from the saved GetDP study, not from illustrative plots.
The report exporter checks the recorded SHA-256 hashes against the current FE
inputs before typesetting. The three supplied TA \texttt{.pro} files are preserved.

\begingroup\small\setlength{\parskip}{7pt}
[1] KU Leuven, \textit{NMES Assignment 2026}, Section 2 and Appendices A--B.
Governing equations, peak-phasor conventions and coaxial reference formula.
Local assignment supplied with the lab.

[2] NKT, \textit{33 kV single core cable XLPE-AL Halogen free}, printed pp.~19--20,
\(1\times240+50~\mathrm{mm^2}\) entry. Dimensions, construction, 33 kV rating and
maximum capacitance \(0.24~\mu\mathrm F/\mathrm{km}\).
\source{https://nkt.widen.net/content/d8gut55fhb/pdf/33kV_1-core_XLPE-AL-RMP-FB-ST_CU_screen_DSENAPP_MV_DS_EN.pdf?u=gj0n1y}{Manufacturer datasheet}.

[3] Authority for Public Services Regulation, Oman, \textit{OES 2: Cables and
Accessories}, Fourth Edition, 2024, 33 kV specification, printed pp.~20--21.
Reference screen thicknesses and resistivity limits used in Section 1;
not an assertion of NKT-specific compound data.
\source{https://www.dcrp-oman.com/documents/DCRP\%20Updates/OES2\%204th\%20Edition-Circular.pdf}{Official utility/regulator document}.

[4] Lee et al., \textit{Comparative Analysis of XLPE and Thermoplastic
Insulation-Based HVDC Power Cables}, Energies 16(1), 167 (2023), Table 3.
Adopted semiconducting-PE permittivity and conductivity reference.
\source{https://doi.org/10.3390/en16010167}{doi:10.3390/en16010167}.

[5] R.~Chippendale, University of Southampton, \textit{Parameters for extra large
cables}, report 15242-RE2-v7, 14 October 2016, Tables 3--4.
Adopted XLPE permittivity.
\source{https://www.transmissionexcel.com/wp-content/uploads/2016/11/Parameters-for-extra-large-cables-FINAL.pdf}{Report}.

[6] COMSOL, \textit{Submarine Cable 8: Inductive Effects 3D}, material discussion.
Adopted XLPE conductivity; the example's cable geometry is not used.
\source{https://doc.comsol.com/6.4/doc/com.comsol.help.models.acdc.submarine_cable_08_inductive_effects_3d/submarine_cable_08_inductive_effects_3d.html}{Official technical documentation}.

[7] Eskom, \textit{Specification for XLPE-insulated power cables and accessories
for systems with nominal voltages of 44 kV to 132 kV}, 240-56030625, Rev.~3,
4 July 2022, \S3.3(i), p.~7; Annex B items 3.29--3.30, p.~22.
External RMS screen-stress comparisons; applicability limitation in Q2.4.
\source{https://www.etenders.gov.za/home/Download/?blobName=b55b97a6-fb2e-492f-8eb0-01bc9574c555.pdf\&downloadedFileName=240-56030625_Specificatin+for+HV+XLPE+cables+and+accessories.pdf}{Official procurement copy}.

[8] P.~Dular and C.~Geuzaine, \textit{GetDP Reference Manual}; C.~Geuzaine and
J.-F.~Remacle, \textit{Gmsh Reference Manual}.
\source{https://getdp.info/doc/texinfo/getdp.html}{GetDP} and
\source{https://gmsh.info/doc/texinfo/gmsh.html}{Gmsh} official documentation.
Complex nodal formulation, physical groups, mesh formats and post-processing.
\endgroup

\subhead{How to reproduce and inspect the results}
From \path{Source-code/section_2}, run
\begin{verbatim}
python3 run.py --study
python3 run.py --open
\end{verbatim}
Open the generated \path{results/healthy/view.geo} or
\path{results/void/view.geo} in Gmsh. Field, potential and current-density
views are available. The central unmeshed disks are the ideal aluminium electrodes;
visible outer gaps include omitted copper/PE regions, not physical air spaces.

\path{Source-code/README.md} explains the complete codebase.
\path{section_2/README.md} gives the detailed model assumptions;
\path{validation/summary.json} records ten cases, checks and tool versions.
The PDF is rebuilt with \path{LaTeX/build_section_2.sh}.
Codex assisted with coding, source lookup, verification and this LaTeX report;
this use is declared in the README as required by the assignment.
\end{document}
